% chapters/05_diskussion.tex

\chapter{Diskussion}
\label{chap:diskussion}

Die Ergebnisse aus \cref{chap:ergebnisse} wirken auf den ersten Blick
widersprüchlich: ein sauber konvergiertes Training, aber eine Erfolgsrate von
$0/20$ im geschlossenen Regelkreis. Dieses Kapitel zeigt, dass dieser Befund
keine Anomalie, sondern ein in der Literatur dokumentiertes und erwartbares
Ergebnis ist, und diskutiert seine Konsequenzen für die Bewertung der
Modellqualität.

\section{Einordnung des Null-Resultats}
\label{sec:einordnung-null}

Die zentrale Erkenntnis lautet: \textbf{$0\,\%$ Closed-Loop-Erfolg in einer
nicht-photorealistischen Simulation mit eingefrorenem Realbild-Encoder ist das
erwartete Ergebnis des Real$\to$Sim-Gaps -- kein Urteil über das Modell.} Die
kanonische Referenz hierfür ist SIMPLER~\cite{simpler2024}: Dort erreicht etwa die
Octo-Politik in der Aufgabe \enquote{Pick Coke Can} \textbf{$0\,\%$} Erfolg,
solange das simulierte Erscheinungsbild nicht an die realen Bilder angeglichen
wird, und springt nach \emph{Visual Matching} auf~$29{,}3\,\%$. Die in dieser
Arbeit verwendete Szene -- einfache, mattierte Quader vor sauberem Hintergrund --
ist für einen auf Realbildern trainierten Encoder geradezu der ungünstigste Fall.

Verstärkt wird dies durch die in \cref{sec:domain-gap-ergebnis} gemessene
Verteilung des Domain-Gaps: Nicht der Durchschnitt ist das Problem, sondern der
\emph{eine} kritische Kameraeingang (\texttt{cam\_left\_wrist}, Distanz~0,427
in der Erstmessung; nach Kamera-Neukalibrierung und Albedo-Fixes noch~0,356,
\cref{sec:neukalibrierung}).
Da das \ac{vlm} alle vier Kameras gemeinsam einbettet, propagiert ein einzelner
stark abweichender Eingang in das gesamte Vision-Language-Embedding -- die in
\cref{sec:imitation-learning} beschriebenen \emph{compounding errors} tun dann
ihr Übriges, bis die Politik in den beobachteten \enquote{eingefrorenen} Zustand
kollabiert.

Den Schritt von der \emph{Bildähnlichkeit} zur \emph{Verhaltensfolge} leistet
das \texttt{span}-Gate (\cref{sec:span-gate}). Es beantwortet die Frage, die
eine Embedding-Distanz allein offenlässt: ob der gemessene Bildunterschied
tatsächlich das Verhalten ändert. Dieselbe Politik kommandiert auf echten
Datensatzbildern die volle Greifbewegung (Verhältnis~1,00), in der Simulation
nur ein Fünftel davon. Damit ist die naheliegendste Gegenhypothese -- die
Politik habe das Greifen schlicht nie gelernt -- ausgeschlossen, und zwar
durch eine Messung, deren beide Seiten von derselben Politik stammen und sich
allein in der Beobachtung unterscheiden. Der in
\cref{sec:vorbehalt-greifzahlen} formulierte Vorbehalt schränkt dabei die
simulationsseitige Hälfte dieser Messung ein, nicht die reale.

\begin{infobox}[Die richtige Metrik für dieses Modell]
  NVIDIAs offizielle Checkpoint-Evaluation ist nicht die Closed-Loop-Sim, sondern
  der \textbf{Open-Loop-Action-\ac{mse} auf gehaltenen echten Trajektorien}
  (\cref{sec:eval-methodik}). Die in \cref{sec:diagnose} dokumentierte gute
  Open-Loop-Armvorhersage spricht dafür, dass das Modell die Aufgabe im Prinzip
  gelernt hat; das Closed-Loop-Versagen ist primär ein \emph{Wahrnehmungs}-, kein
  \emph{Politik}-Problem.
\end{infobox}

\section{Eingefrorener vs.\ mittrainierter Vision-Encoder}
\label{sec:frozen-vs-tuned}

Ein naheliegender Reflex ist, den Vision-Encoder mitzutrainieren
(\texttt{TUNE\_VISUAL=1}), um ihn an die Simulation anzupassen. Dieser Schluss ist
jedoch trügerisch: Domain-Invarianz ist eine \emph{im Training erlernte}
Eigenschaft, kein zur Evaluationszeit umlegbarer Schalter. Ein Encoder, der
ausschließlich auf Realbildern feinabgestimmt wird, wird dadurch \emph{nicht}
robuster gegenüber Isaac-Sim-Renderings -- er spezialisiert sich im Gegenteil noch
stärker auf das exakte Erscheinungsbild der Trainingsbilder. Ein
\texttt{TUNE\_VISUAL}-Training hilft gegen den Domain-Gap nur dann, wenn der Encoder
während des Trainings tatsächlich Simulationsbilder oder durch \acf{dr} variierte
Bilder sieht~\cite{tobin2017domain}.

Hinzu kommt ein praktisches Risiko: Bei nur~$\approx$~301~Demonstrationen droht
beim Auftauen des 2-Mrd.-Parameter-Encoders \emph{Catastrophic Forgetting} der
generischen, vortrainierten Merkmale -- gerade jener Merkmale, die einen
eingefrorenen Encoder vergleichsweise sim-tolerant machen. Frozen-Encoder sind
für die Sim-Evaluation also tendenziell robuster, getunte Encoder besser für den
Real-zu-Real-Transfer. Diese Abwägung ist in \doc{docs/ergebnisse/sim-bewertung.md}
und \doc{docs/ergebnisse/domain-gap-analyse.md} weiter ausgeführt.

Diese Argumentation ist kein bloßes Gedankenexperiment: Der in
\cref{sec:vision-lauf} geschilderte zweite Lauf hat \texttt{TUNE\_VISUAL=1} genau
unter diesen Bedingungen erprobt -- und das vorhergesagte Verhalten reproduziert.
Bei identischem Trainingsbefund (Endloss~$\approx$~0,009) verschlechterte sich das
Closed-Loop-Verhalten von \enquote{Anfahren-dann-Einfrieren} zu einem reinen
Politik-Kollaps (nur noch Arm-Rückzug). Das Mittrainieren des Encoders auf reinen
Realdaten ist damit als Sackgasse \emph{empirisch belegt} und nicht nur theoretisch
plausibilisiert; sinnvoll wäre es erst zusammen mit sim- oder
\acs{dr}-variierten Trainingsbildern.

\section{Validität der Evaluationsmethodik}
\label{sec:validitaet}

Die kritische Selbstprüfung der Sim-Eval kommt zu einem klaren Urteil: Die
Simulation ist \textbf{technisch korrekt gebaut, aber als Bewertungsinstrument
für dieses Modell methodisch am Ziel vorbei.} Der Steuerungspfad ist durch den
Replay-Test nachweislich korrekt (\cref{sec:kalibrierung}), und die
Inferenz-Hälfte reproduziert einen unabhängigen, publizierten Benchmark auf
0,1~Prozentpunkte genau (\cref{sec:robocasa}); die Closed-Loop-Erfolgsrate
misst jedoch primär den Domain-Gap, nicht die Modellgüte. Bei der
Durchsicht der Eval-Implementierung wurden zudem einige kleinere, teils
kosmetische Mängel identifiziert (\cref{tab:codebefunde}), die für die
Interpretation relevant sind, das Null-Resultat aber \emph{nicht} erklären.

\begin{table}[htb]
  \centering
  \begin{tabularx}{\linewidth}{lX}
    \toprule
    \textbf{Priorität} & \textbf{Befund} \\
    \midrule
    hoch    & Irreführender Kommentar \enquote{Arme = relativ}, obwohl der Code
              serverseitig alle 28~Dimensionen absolut anwendet. \\
    mittel  & \texttt{dr\_enabled=True} zur Evaluationszeit ist für einen
              eingefrorenen Encoder wirkungslos (\ac{dr} ist eine
              \emph{Trainings}-Technik) und sollte deaktiviert sein. \\
    mittel  & Die Steuerfrequenz beträgt real~28,57~Hz statt der im Label
              genannten~30~Hz -- für absolute Positionsziele unkritisch. \\
    niedrig & Die Erfolgsbedingung (vollständige Dreier-Säule) ist gegen die
              tatsächliche Datensatz-Task abzugleichen. \\
    \bottomrule
  \end{tabularx}
  \caption{Bei der Methodik-Prüfung identifizierte Code-Befunde
    (\doc{docs/ergebnisse/sim-bewertung.md}). Keiner ist die Ursache des
    Null-Resultats.}
  \label{tab:codebefunde}
\end{table}

Mehrere zunächst verdächtigte Faktoren konnten durch Recherche \emph{entlastet}
werden: der ausgeführte Aktionshorizont (8~von~16, eine Lehrbuch-Konvention für
Receding-Horizon-Politiken), die geringe Frequenzabweichung und die
absolut/relativ-Umsetzung der Aktionen. Sie sind nicht die Fehlerursache.

\section{Messartefakte und ihre Lehren}
\label{sec:messartefakte}

Die Diagnose-Serien vom August~2026
(\cref{sec:neukalibrierung,sec:grasp-befund,sec:geometrie}) förderten über die
Sachbefunde hinaus vier wiederkehrende \emph{Messartefakt-Muster} zutage, die
für simulationsgestützte Robotik-Evaluationen verallgemeinerbar sind:

\begin{description}
  \item[Zirkuläre Prüfgrößen.] Der Kamera-Orientierungsfehler
    (\cref{sec:neukalibrierung}) blieb fünf Diagnose-Läufe lang unentdeckt,
    weil die Kontrollgröße aus \emph{derselben} fehlerhaften
    Koordinaten-Umrechnung stammte wie die geprüfte Kamerapose: Eine Größe,
    die auf Hin- und Rückweg durch dieselbe Transformation läuft, kann diese
    Transformation nicht prüfen. Erst eine unabhängige Messquelle -- der
    direkte Zugriff auf die USD-Szene -- machte die Abweichung von bis
    zu~137° sichtbar. Kontrollmessungen müssen den geprüften Codepfad
    umgehen.
  \item[Körper-Frames sind keine Kontaktflächen.] Der Abstand
    \enquote{Hand--Würfel} wurde in der Greif-Diagnostik dreimal
    nacheinander von einem falschen Referenzpunkt aus gemessen -- Handwurzel,
    Handfläche, distales Fingergelenk --, dessen Ursprung jeweils
    zentimeterweit (zuletzt 5,2~cm) vor der tatsächlichen Kontaktfläche der
    Fingerkuppe liegt. Jede dieser Messungen erzeugte eine falsche
    Zwischen-Schlussfolgerung; so erwies sich der zunächst alarmierende
    Befund \enquote{die Arme nähern sich dem Würfel nie} als reines
    Messartefakt der Handwurzel-Referenz. Bei einem Würfel von 5~cm
    Kantenlänge entscheidet die Wahl des Referenzpunkts über das Vorzeichen
    des Ergebnisses.
  \item[Stille Sättigung statt Fehlermeldung.] Der widerrufene
    Vorzeichen-Fix (\cref{sec:vorzeichen-widerruf}) blieb monatelang
    unentdeckt, weil die Simulation Gelenkziele außerhalb ihrer Grenzen nicht
    zurückweist, sondern stumm auf den Grenzwert klemmt. Aus Sicht jeder
    darüberliegenden Auswertung lieferte das System weiterhin plausible
    Zahlen -- ein stillstehendes Gelenk ist von einem Gelenk, das seine
    Zielposition erreicht hat, im Protokoll nicht zu unterscheiden. Wo ein
    Stellglied sättigen kann, muss der Sättigungszustand selbst protokolliert
    werden, sonst wird ein Defekt als Messwert gelesen.
  \item[Prüfdaten sind nicht Nutzdaten.] Die ursprüngliche Begründung für den
    Wechsel des Gierwinkel-Schätzers (\cref{sec:gierwinkel}) musste
    zurückgezogen werden: Sie war auf \emph{annotierten Debug-Bildern} statt
    auf echten Datensatzbildern erhoben worden und maß damit die eingezeichneten
    Hilfslinien mit. Aufgefallen ist das an exakt gesättigten Farbwerten
    ($255,0,0$ und $255,255,255$), wie sie in Fotografien nicht vorkommen.
    Auffällig \enquote{saubere} Werte sind daher weniger ein Qualitätsindiz als
    ein Anlass zu prüfen, ob überhaupt die richtigen Daten gemessen wurden.
\end{description}

Beide Muster teilen eine Konsequenz: Quantitative Diagnostik ist nur so
belastbar wie die Unabhängigkeit und die physikalische Bedeutung ihrer
Messpunkte. Die im Projekt daraus abgeleiteten Regeln -- Kontrollgrößen aus
unabhängigen Quellen beziehen, Abstände an Kontaktflächen statt an
Körper-Frames messen und jeden Messwert mit seinem Bezugsrahmen protokollieren
-- sind in \doc{docs/weiterfuehrend/rl-anleitung.md} operationalisiert.

\section{Limitierungen}
\label{sec:limitierungen}

\subsection{Vorbehalt gegen alle simulationsseitigen Greifzahlen}
\label{sec:vorbehalt-greifzahlen}

Die gewichtigste Einschränkung dieser Arbeit ergab sich erst ganz am Ende und
betrifft rückwirkend einen Großteil ihrer Messwerte. Wie in
\cref{sec:vorzeichen-widerruf} geschildert, war die für die Fingergrundgelenke
vorgenommene Vorzeichen-Spiegelung selbst fehlerhaft: Sie drehte vier Gelenke
aus ihrem zulässigen Bereich, wo die Simulation sie stumm auf den Grenzwert
klemmte. Da die Spiegelung sowohl auf die \emph{Beobachtung} als auch auf die
\emph{Aktion} wirkte, betraf sie jeden Laufmodus. Konkret heißt das: Sämtliche
in dieser Arbeit berichteten simulationsseitigen Greifgrößen -- die
Erfolgsrate~$0/20$, die Fingerspannen von 19\,\%, 20,5\,\% und 27,6\,\% -- sind
mit einer Hand entstanden, die sich nicht vollständig schließen konnte.

Für die Belastbarkeit der Argumentation ist entscheidend, sauber zu trennen,
was dieser Vorbehalt berührt und was nicht:

\begin{description}
  \item[Nicht betroffen.] Die SigLIP-Messung des Domain-Gaps
    (\cref{sec:domain-gap-ergebnis,sec:neukalibrierung}) ist rein bildbasiert
    und kennt keine Gelenkkonventionen; sie belegt unabhängig, dass ein
    visueller Abstand besteht. Ebenso unberührt ist die externe
    Pipeline-Validierung (\cref{sec:robocasa}), die auf einem anderen
    Embodiment in einem anderen Simulator läuft und die DEX3-Hand gar nicht
    verwendet. Auch die \emph{reale} Seite des \texttt{span}-Gates
    (Verhältnis~1,00) ist sauber: Sie wird open-loop auf echten Bildern ohne
    jede Simulationsbeteiligung erhoben.
  \item[Betroffen.] Alle in der Simulation erhobenen Greif- und
    Fingerspannenwerte. Für sie existiert seit dem Widerruf eine
    \emph{konkurrierende Erklärung}: Die Politik sah in der Simulation nicht
    nur fremde Bilder, sondern zusätzlich einen verfälschten propriozeptiven
    Zustand -- und beides kann eine zu geringe kommandierte Fingerspanne
    erzeugen.
\end{description}

Die qualitative Kernaussage der Arbeit übersteht das: Dass die Politik das
Greifen \emph{gelernt} hat, zeigt die unbelastete reale Seite des
\texttt{span}-Gates, und dass ein erheblicher visueller Abstand \emph{besteht},
zeigt die unbelastete SigLIP-Messung. Was der Vorbehalt schwächt, ist die
\emph{Alleinzuschreibung} des Sim-Versagens an den visuellen Domain-Gap: Ein
propriozeptiver Beobachtungsfehler steht nun als zweite, nicht abgetrennte
Ursache daneben. Die saubere Konsequenz wäre, die simulationsseitigen
Messungen mit korrigierter Geometrie zu wiederholen; das war zum
Redaktionsschluss nicht mehr möglich. Bis dahin sind die genannten Prozentwerte
als \emph{obere Schranken für das Ausmaß des visuellen Effekts} zu lesen, nicht
als dessen Messung.

\subsection{Datenmenge und Datenqualität}
Mit~$\approx$~301~Demonstrationen liegt der Datensatz am unteren Rand des für
\ac{vla}-Fine-Tuning Üblichen (typischerweise 500--2000). In Kombination mit
reinem \acl{bc} begünstigt das die in \cref{sec:imitation-learning} beschriebenen
\emph{compounding errors}. Zudem waren die Fingerdimensionen schon in der
Open-Loop-Vorhersage verrauscht -- ein Hinweis auf grenzwertige Datenqualität
oder -menge gerade für die feinmotorischen DEX3-Gelenke.

\subsection{Fehlende Validierung und blinde Checkpoint-Auswahl}
Da während des Trainings nicht validiert wurde, fehlte ein Signal zur Auswahl
des besten Checkpoints; Schritt~175\,000 wurde mangels Alternative genommen,
obwohl die Loss-Kurve bereits ab~$\approx$~100\,000 plateaut. Der zusätzlich
aufgedeckte inaktive Train/Test-Split (\cref{sec:split-befund}) bedeutet, dass
für diesen Lauf keine sauber gehaltene Testmenge existiert. Die Einschränkung
ist dabei nicht selbst verschuldet, sondern strukturell: Der verwendete Fork
erzwingt \texttt{eval\_strategy="no"} per Assertion.

Diese Limitierung ist inzwischen \emph{quantifiziert} und behoben. Der dritte
Lauf (\cref{sec:lauf3}) verfügte über eine echte Testmenge, und der
nachgelagerte Checkpoint-Sweep zeigte eine U-Kurve: Der letzte Checkpoint ist
um 25\,\% schlechter als der beste. Damit ist das Overfitting im Projekt
erstmals gemessen statt vermutet -- und rückwirkend belegt, dass die Läufe~1
und~2 mit hoher Wahrscheinlichkeit den falschen Checkpoint ausgewertet haben.
Wie stark dieser Effekt die Closed-Loop-Befunde beeinflusst hat, lässt sich
nicht mehr rekonstruieren; die Richtung ist jedoch klar, und für künftige Läufe
ist die Auswahl durch den Sweep methodisch abgesichert.

\subsection{Abhängigkeit von der Simulator-Version}
Der Wechsel von Isaac~Sim~4.x auf~6.0 (\cref{sec:neukalibrierung}) zeigte, dass
zentrale Validierungsergebnisse an die Simulator-Version gebunden sind: Die im
Juni belegte Greif-Physik (2,8~cm Hubhöhe im Replay) ließ sich nach der
Migration zunächst nicht reproduzieren (\cref{sec:grasp-befund}). Die Klärung
ergab am Ende \emph{Entwarnung} für den Simulator -- mit korrekt platziertem
Würfel hebt auch Isaac~Sim~6.0 an (7,9~cm) --, kostete aber mehrere Wochen
Diagnose. Der eigentliche Fehler lag nicht in der neuen Version, sondern in der
Szenengeometrie (\cref{sec:geometrie}). Genau das ist die verallgemeinerbare
Lehre: Ein Versionswechsel verschiebt die Aufmerksamkeit auf den Simulator und
verdeckt damit Fehler, die schon vorher bestanden. Simulationsgestützte Befunde
sollten stets mit der Simulator-Version ausgewiesen und nach Versionswechseln
erneut geprüft werden -- die Prüfung sollte jedoch die Szene einschließen, nicht
nur den Simulator.

\subsection{Sim-to-Real-Transfer}
Die Evaluation fand ausschließlich in Simulation statt. Der entscheidende Test --
das Verhalten auf dem \emph{physischen} G1 -- steht aus. Gerade weil das
Closed-Loop-Versagen ein Sim-spezifisches Wahrnehmungsproblem ist, ist plausibel,
dass das Modell auf realer Hardware (mit realen Kamerabildern, die der
Trainingsverteilung entsprechen) deutlich besser abschneidet. Der reale Roboter
bleibt damit der eigentliche Goldstandard der Bewertung.

\section{Methodische Konsequenz}
\label{sec:konsequenz}

Aus alledem folgt eine klare Empfehlung für die weitere Bewertung dieses Modells,
die der in der Literatur etablierten Hierarchie entspricht:

\begin{enumerate}
  \item \textbf{Open-Loop-Action-\ac{mse}} als primäre, sofort verfügbare
    Metrik für die Modellqualität (kein Domain-Gap).
  \item \textbf{Reale Roboter-Rollouts} als Goldstandard.
  \item \textbf{Closed-Loop-Sim} nur mit SIMPLER-artigem Visual
    Matching~\cite{simpler2024} oder nach Aufnahme von Sim-Bildern in den
    Trainingssatz -- andernfalls taugt die Sim als qualitatives Demo- und
    Debug-Werkzeug, nicht als Modell-Benchmark.
\end{enumerate}

Die Simulation hat in dieser Arbeit ihren Zweck dennoch erfüllt: Sie hat -- gerade
durch ihr \enquote{Scheitern} -- den Domain-Gap als zentrale Herausforderung
sichtbar gemacht und damit den Weg für die in \cref{chap:weiterfuehrend}
skizzierten Folgearbeiten gewiesen.
